\documentclass[a4paper,11pt]{article}

\usepackage[utf8]{inputenc}
\usepackage[T1]{fontenc}
\usepackage{lmodern}
\usepackage[a4paper,margin=2.2cm]{geometry}
\usepackage{amsmath,amssymb,amsfonts}
\usepackage{booktabs}
\usepackage{tabularx}
\usepackage{array}
\usepackage{makecell}
\usepackage{xcolor}
\usepackage{microtype}
\usepackage{hyperref}
\usepackage{fancyhdr}

\hypersetup{
    colorlinks=true,
    linkcolor=blue!70!black,
    citecolor=blue!70!black,
    urlcolor=blue!70!black
}

\pagestyle{fancy}
\fancyhf{}
\rhead{\small Formal Verification of Leaky Assertions in TLA+/TLC}
\lhead{\small Experimental Evaluation Report}
\cfoot{\thepage}

\title{\textbf{Formal Verification of Information Flow and Side-Channel Leaks via Leaky Assertions in TLA+/TLC}\\[0.5ex]
\large Experimental Evaluation Report and Model Analysis}
\author{\textbf{Formal Methods and Security Research Group}}
\date{\today}

\begin{document}

\maketitle


\section{Experimental Evaluation Summary}

Table~\ref{tab:benchmarks} summarizes the model checking statistics obtained from TLC runs across all six benchmark suites. The state-space exploration metrics reflect breadth-first model checking on each specification.

\begin{table}[htbp]
\centering
\small
\renewcommand{\arraystretch}{1.25}
\begin{tabularx}{\textwidth}{@{} >{\raggedright\arraybackslash}X ccccc @{}}
\toprule
\textbf{Benchmark Program} & 
\makecell{\textbf{No. of Threads} \\ \textbf{($n$)}} & 
\makecell{\textbf{Total} \\ \textbf{States}} & 
\textbf{Depth} & 
\makecell{\textbf{TLC Time} \\ \textbf{(ms)}} & 
\makecell{\textbf{Leak Found?} \\ \textbf{(Yes/No)}} \\
\midrule
\texttt{ConcurrentLeak}            & 2 & 44  & 5  & 571 & \textbf{Yes} \\
\texttt{InternalTiming}            & 2 & 20  & 10 & 706 & \textbf{Yes} \\
\texttt{InternalTimingTransformed} & 2 & 22  & 10 & 586 & \textbf{No}  \\
\texttt{ModSqLeaks}                & 1 & 322 & 21 & 466 & \textbf{Yes} \\
\texttt{WebServiceLeak}            & 1 & 17  & 5  & 485 & \textbf{Yes} \\
\texttt{SeqLock}                   & 3 & 3   & 2  & 729 & \textbf{Yes} \\
\bottomrule
\end{tabularx}
\caption{Model checking results across all benchmark programs evaluated in TLC.}
\label{tab:benchmarks}
\end{table}

\noindent\textbf{Note on Transformed and Mitigated Programs:}
\begin{itemize}
    \item In \texttt{InternalTimingTransformed}, the branch-padding mitigation successfully eliminates execution-time variance ($\Delta = 9$ across all inputs), resulting in \textbf{No} leak found (all non-interference invariants pass).
    \item For \texttt{WebServiceLeak} and \texttt{SeqLock}, the table presents the vulnerable baseline configuration (\texttt{Fixed = FALSE}). When the mitigation toggle is engaged (\texttt{Fixed = TRUE}), the leak is completely eliminated, and TLC reports 0 violations across all reachable states.
\end{itemize}

\bigskip

\section{Benchmark Program Writeups}

\subsection{ConcurrentLeak (\texttt{ConcurrentLeaks/ConcurrentLeak.tla})}

\subsubsection{1. How the Program is Leaking Information}
The program models two concurrent processes, $T_1$ and $T_2$, synchronizing via a shared binary semaphore (\texttt{sem}) and sharing variable $v$. Process $T_1$ unconditionally enters a critical section, emits public output event \texttt{"a"}, increments $v$, emits public event \texttt{"b"}, and releases the semaphore. Process $T_2$ emits public entry event \texttt{"c"}, records auxiliary entry timestamp $T_{\text{start}} := t$, and then conditionally executes an internal critical section based on the secret:
\begin{itemize}
    \item When $h = 1$: $T_2$ contends for the semaphore, waits for $T_1$ to finish, increments $v := v + 2$, and releases the semaphore before emitting exit event \texttt{"d"}.
    \item When $h = 0$: $T_2$ skips the critical section completely and immediately emits exit event \texttt{"d"}.
\end{itemize}
This introduces two distinct side-channel leaks:
\begin{enumerate}
    \item \textbf{Timing Side-Channel}: When $h = 0$, $T_2$ exits fast ($\Delta \le 3$). When $h = 1$, semaphore contention forces $T_2$ to serialize behind $T_1$, yielding prolonged elapsed execution time ($\Delta \ge 7$).
    \item \textbf{Order Covert Channel}: Interleaving produces observable output event traces in sequence \texttt{out}. When $h = 1$, serialization forces the trace prefix \texttt{"acdb"}, whereas when $h = 0$, $T_2$ can finish earlier producing interleavings like \texttt{"cadb"} or \texttt{"cdab"}.
\end{enumerate}

\subsubsection{2. Detail About the Secret Variables}
\begin{itemize}
    \item \textbf{Variable}: $h \in \{0, 1\}$.
    \item \textbf{Classification}: Secret (high confidentiality).
    \item \textbf{Role}: Directly governs whether $T_2$ attempts to acquire the shared semaphore and access the shared variable $v$.
\end{itemize}

\subsubsection{3. Leaky Assertions Used to Identify Leaks in TLC}
TLC verifies the following invariants at the exit public statement $l_7$:
\begin{itemize}
    \item \textbf{Soundness of Timing Leaky Assertion (Evaluates to TRUE)}:
    \begin{align*}
        \text{LA\_Timing} &\equiv \text{Observed} \implies ((\Delta \ge \Delta_{\text{hi}}) \iff (h = 1))\\
        \text{LA\_Time\_Low} &\equiv (\text{Observed} \land \Delta \le \Delta_{\text{lo}}) \implies (h = 0)\\
        \text{LA\_Time\_High} &\equiv (\text{Observed} \land \Delta \ge \Delta_{\text{hi}}) \implies (h = 1)
    \end{align*}
    where $\Delta_{\text{lo}} = 3$ and $\Delta_{\text{hi}} = 7$.
    \item \textbf{Soundness of Order Leaky Assertion (Evaluates to TRUE)}:
    $$\text{LA\_Trace1} \equiv (\text{Observed} \land \text{out} = \langle \text{"a"}, \text{"c"}, \text{"d"}, \text{"b"} \rangle) \implies (h = 1)$$
    \item \textbf{Leak Discovery Invariant (Violated in TLC)}:
    $$\text{NoTimingLeak} \equiv \text{Observed} \implies (\Delta = \Delta_{\text{lo}})$$
    TLC refutes this property, outputting the counterexample trace where $h = 1$ causes $\Delta = 7$.
\end{itemize}

\bigskip

\subsection{InternalTiming (\texttt{InternalTiming/InternalTiming.tla})}

\subsubsection{1. How the Program is Leaking Information}
This model captures the subtle interaction between secret-dependent branching and a deterministic round-robin scheduler. Two threads, $T_1$ and $T_2$, alternate execution turns:
\begin{itemize}
    \item Thread 1 conditionally branches on secret $h$:
    \begin{itemize}
        \item If $h = 1$: executes through labels $l_1 \to l_2 \to l_3 \to l_5$, requiring 3 scheduling turns before reaching public assignment $l := 1$ (\texttt{out := Append(out, "1")}).
        \item If $h = 0$: executes through labels $l_1 \to l_4 \to l_5$, requiring only 2 scheduling turns before reaching $l_5$.
    \end{itemize}
    \item Thread 2 executes sequentially through $m_1 \to m_2 \to m_3$ where it sets $l := 0$ (\texttt{out := Append(out, "0")}).
\end{itemize}
Because Thread 1's branch imbalance alters the scheduler turn allocation, the total global execution time across both threads differs: $\Delta = 8$ clock cycles when $h = 0$, and $\Delta = 9$ clock cycles when $h = 1$. An external observer measuring elapsed time between the public entry statement \texttt{<<"START">>} and public termination \texttt{<<"END">>} directly reconstructs the private branch choice.

\subsubsection{2. Detail About the Secret Variables}
\begin{itemize}
    \item \textbf{Variables}: $h \in \{0, 1\}$ (primary branch secret), $h_2 \in \mathbb{N}$ (cryptographic accumulator updated via $h_2 := 2 \cdot h_2 + l$).
    \item \textbf{Classification}: Secret (high data).
    \item \textbf{Role}: Governs the execution path length in Thread 1, directly perturbing round-robin turn progression.
\end{itemize}

\subsubsection{3. Leaky Assertions Used to Identify Leaks in TLC}
Evaluated at process termination (\texttt{Observed}):
\begin{itemize}
    \item \textbf{Attacker Deductive Leaky Assertions (PASS in TLC)}:
    \begin{align*}
        \text{Calib0} &\equiv (\text{Observed} \land h = 0) \implies \Delta = 8\\
        \text{Calib1} &\equiv (\text{Observed} \land h = 1) \implies \Delta = 9\\
        \text{LA\_Timing} &\equiv \text{Observed} \implies ((\Delta \ge 9) \iff (h = 1))
    \end{align*}
    \item \textbf{Security Invariant (FAILS in TLC)}:
    $$\text{NoTimingLeak} \equiv \text{Observed} \implies (\Delta = 8)$$
    TLC finds an execution path with $h = 1$ yielding $\Delta = 9$, disproving timing non-interference.
\end{itemize}

\bigskip

\subsection{InternalTimingTransformed (\texttt{InternalTimingTransformed/Transformed.tla})}

\subsubsection{1. How the Program is Leaking Information (Mitigation Analysis)}
\texttt{InternalTimingTransformed} represents the mitigated countermeasure specification. To eliminate the internal timing covert channel caused by the branch imbalance in Thread 1, source-level \emph{branch padding} is introduced.
\begin{itemize}
    \item On the shorter path ($h = 0$), an auxiliary padding step is inserted:
    $$\texttt{l4\_pad: await turn = 1; skip; t := t + 1; turn := 2; goto l5;}$$
    \item This balances the turn consumption: both branches now consume exactly 3 scheduler turns before reaching $l_5$.
    \item As a result, the total elapsed execution time becomes strictly invariant: $\Delta = 9$ across both $h = 0$ and $h = 1$. The side-channel leak is eliminated.
\end{itemize}

\subsubsection{2. Detail About the Secret Variables}
\begin{itemize}
    \item \textbf{Variables}: $h \in \{0, 1\}$ and $h_2 \in \mathbb{N}$.
    \item \textbf{Classification}: Secret (high data).
    \item \textbf{Role}: Controls branching in Thread 1; however, both branches are now structurally balanced in scheduling turns and execution cycles.
\end{itemize}

\subsubsection{3. Leaky Assertions Used to Identify Leaks in TLC}
\begin{itemize}
    \item \textbf{Security Invariant (PASSES in TLC)}:
    $$\text{ConstantTime} \equiv \text{Observed} \implies (\Delta = 9)$$
    $$\text{NoTimingLeak} \equiv \text{Observed} \implies (\Delta = 9)$$
    TLC checks all 20 reachable states and confirms that execution time is completely independent of $h$.
    \item \textbf{Attacker Leaky Assertion Defeat}:
    The previous attacker leaky assertion $\text{LA\_Timing} \equiv \text{Observed} \implies (\Delta \ge 9 \iff h = 1)$ is invalidated, proving the attacker can no longer distinguish $h = 0$ from $h = 1$.
\end{itemize}

\bigskip

\subsection{ModSqLeaks (\texttt{ModSqLeaks/ModSqLeaks.tla})}

\subsubsection{1. How the Program is Leaking Information}
This specification models the classic square-and-multiply algorithm used in modular exponentiation ($y^K \pmod N$), foundational to RSA and Diffie-Hellman implementations.
\begin{itemize}
    \item For each bit $i$ of the exponent $K$:
    \begin{itemize}
        \item The algorithm unconditionally executes a modular squaring operation ($t := t + 1$).
        \item If the secret bit $K[i] = 1$, it additionally executes a modular multiplication ($t := t + 1$).
        \item If $K[i] = 0$, the multiplication is bypassed (when $\texttt{Fixed = FALSE}$).
    \end{itemize}
    \item \textbf{The Leak}: Total execution time $\Delta = \text{Len}(K) + \text{HW}(K)$, where $\text{HW}(K)$ is the Hamming weight (number of 1-bits) of the secret exponent. An observer measuring execution time learns the exact Hamming weight or bit threshold of the private key.
    \item \textbf{Mitigation ($\texttt{Fixed = TRUE}$)}: A dummy multiplication is performed when $K[i] = 0$, guaranteeing strict constant-time execution ($\Delta = 2 \cdot \text{Len}(K)$).
\end{itemize}

\subsubsection{2. Detail About the Secret Variables}
\begin{itemize}
    \item \textbf{Variable}: $K \in \text{Seq}(\{0, 1\})$, representing the binary representation of the cryptographic key.
    \item \textbf{Derived Secret}: $\text{HW}(K) = \sum_{i} K[i]$ (Hamming weight).
    \item \textbf{Classification}: Highly confidential cryptographic secret.
\end{itemize}

\subsubsection{3. Leaky Assertions Used to Identify Leaks in TLC}
\begin{itemize}
    \item \textbf{Sound Leaky Assertion (Valid in TLC)}:
    $$\text{LeakyAssertion} \equiv \text{Observed} \implies (\Delta \ge \text{Threshold} \iff \text{HW}(K) \ge \text{SecretThreshold})$$
    Formally validates that execution time directly reveals private key bit density.
    \item \textbf{Non-Interference Invariant (FAILS under vulnerable mode)}:
    $$\text{ConstantTime} \equiv \text{Observed} \implies (\Delta = \text{ExpectedConstantTime})$$
    Under $\texttt{Fixed = FALSE}$, TLC violates this invariant and returns traces of varying key weights.
\end{itemize}

\bigskip

\subsection{WebServiceLeak (\texttt{WebServiceLeak/WebServiceLeak.tla})}

\subsubsection{1. How the Program is Leaking Information}
This model captures a user enumeration timing attack against a web authentication endpoint:
\begin{itemize}
    \item A client sends a login request (\texttt{l0: out := Append(out, "REQ")}).
    \item The server queries the database (\texttt{Lookup: t := t + LookupCost}).
    \item \textbf{Vulnerable Path ($\texttt{Fixed = FALSE}$)}:
    \begin{itemize}
        \item If the username does not exist ($\texttt{valid\_user = FALSE}$), the server immediately aborts and returns (\texttt{OutputError: result := "error message"}). Total duration: $\Delta = \text{LookupCost} = 1$.
        \item If the username exists ($\texttt{valid\_user = TRUE}$), but the password is wrong, the server performs an expensive cryptographic password hash (\texttt{HashPassword: t := t + HashCost}) before returning the identical error message. Total duration: $\Delta = \text{LookupCost} + \text{HashCost} = 11$.
    \end{itemize}
    \item Both paths return the identical generic string \texttt{"error message"}, but the response latency differs by an order of magnitude, allowing an external client to deduce whether any target username exists.
    \item \textbf{Mitigation ($\texttt{Fixed = TRUE}$)}: When the username does not exist, the server executes a dummy hash computation (\texttt{DummyHash: t := t + HashCost}), equalizing both failure paths to $\Delta = 11$.
\end{itemize}

\subsubsection{2. Detail About the Secret Variables}
\begin{itemize}
    \item \textbf{Variables}: $\texttt{valid\_user} \in \{\text{TRUE}, \text{FALSE}\}$, $h = \text{IF } \texttt{valid\_user} \text{ THEN } 1 \text{ ELSE } 0$, and $\texttt{password\_match} \in \{\text{TRUE}, \text{FALSE}\}$.
    \item \textbf{Classification}: Secret (confidential database metadata / account existence).
    \item \textbf{Role}: Governs whether the cryptographic hashing routine is bypassed.
\end{itemize}

\subsubsection{3. Leaky Assertions Used to Identify Leaks in TLC}
Evaluated at the HTTP response exit point (\texttt{l\_exit}):
\begin{itemize}
    \item \textbf{Attacker Deductive Leaky Assertions (PASS in TLC)}:
    \begin{align*}
        \text{LA\_Time\_Low} &\equiv (\neg\text{Fixed} \land \text{Observed} \land \text{PublicOut} = \text{"error message"} \land \Delta \le \Delta_{\text{lo}}) \implies (h = 0)\\
        \text{LA\_Time\_High} &\equiv (\neg\text{Fixed} \land \text{Observed} \land \text{PublicOut} = \text{"error message"} \land \Delta \ge \Delta_{\text{hi}}) \implies (h = 1)\\
        \text{LA\_Timing} &\equiv (\neg\text{Fixed} \land \text{Observed} \land \text{PublicOut} = \text{"error message"}) \implies ((\Delta \ge \text{Threshold}) \iff (h = 1))
    \end{align*}
    where $\Delta_{\text{lo}} = 1$, $\Delta_{\text{hi}} = 11$, and $\text{Threshold} = 5$.
    \item \textbf{Security Invariant (FAILS when vulnerable)}:
    $$\text{NoTimingLeak} \equiv (\text{Observed} \land \text{result} = \text{"error message"}) \implies (\Delta = \text{LookupCost} + \text{HashCost})$$
    TLC finds the early abort trace where non-existent users return in $\Delta = 1$ instead of $11$.
\end{itemize}

\bigskip

\subsection{SeqLock (\texttt{SeqLock/SeqLock.tla})}

\subsubsection{1. How the Program is Leaking Information}
Extracted from Graeme Smith's recent work on \emph{Information Flow Security on Persistent Memory} (2026), this benchmark models the Linux reader-writer \texttt{seqlock} mechanism running on Intel-x86 persistent memory (PM) subject to power outages.
\begin{itemize}
    \item Shared memory locations:
    \begin{itemize}
        \item $c$: sequence counter (even = idle, odd = write in progress).
        \item $x_1$: data variable containing either low data or sensitive high secret ($L(x_1) = \text{high}$).
        \item $x_2$: classification flag ($x_2 = 1 \iff x_1$ is high; $L(x_2) = \text{low}$).
    \end{itemize}
    \item \textbf{The Persistency Leak ($\texttt{Fixed = FALSE}$)}:
    Stores to $x_1$, $x_2$, and $c$ reside on different cache lines. Under weak persistency models without explicit cache flushes, stores are committed to persistent memory asynchronously.
    If the store to $x_1$ and the final even store to $c$ persist to PM, while the store to tag $x_2$ lags in volatile buffers when a power outage occurs, rebooting restores an inconsistent state:
    $$c \text{ is even}, \quad x_2 = 0 \text{ (flag indicates low)}, \quad x_1 = 1 \text{ (holds sensitive high secret!)}$$
    A reader executing post-reboot sees $c$ even, reads $x_2 = 0$, concludes the data is low, reads $x_1$, and emits the high secret over the public low channel.
    \item \textbf{Mitigation ($\texttt{Fixed = TRUE}$)}:
    Inserts \texttt{flush c} after the first increment of $c$, and \texttt{flush x1; flush x2} before the final increment of $c$. This ensures $c$ is never persisted as even unless both $x_1$ and $x_2$ are safely committed to persistent memory.
\end{itemize}

\subsubsection{2. Detail About the Secret Variables}
\begin{itemize}
    \item \textbf{Variables}: $h \in \{0, 1\}$, where $h = 1$ denotes high confidential data, and $h = 0$ denotes low public data.
    \item \textbf{Classification}: Secret (high security classification $\Gamma_{x_1} = \text{high}$).
    \item \textbf{Persistent Locations}: \texttt{pm\_x1} (persistent data), \texttt{pm\_x2} (persistent security flag), \texttt{pm\_c} (persistent sequence counter).
\end{itemize}

\subsubsection{3. Leaky Assertions Used to Identify Leaks in TLC}
\begin{itemize}
    \item \textbf{Sound Leaky Assertion (PASSES in TLC)}:
    $$\text{LA\_PersistentLeak} \equiv (\neg\text{Fixed} \land \text{crashed} \land \text{ReaderTerminated} \land (\exists k \in 1..\text{Len}(\text{out}) : \text{out}[k] = 1)) \implies (h = 1)$$
    Proves that an observer observing output value $1$ on the low reader output channel after a crash soundly infers that the high secret was $h = 1$.
    \item \textbf{Value-Dependent Non-Interference Property (FAILS in vulnerable mode)}:
    $$\text{NoInformationLeak} \equiv \forall k \in 1..\text{Len}(\text{out}) : (\text{out}[k] = 1) \implies (h = 0)$$
    When $\texttt{Fixed = FALSE}$, TLC finds the crash-and-reboot trace where the high secret ($h = 1$) is leaked through the reader's low channel, producing the counterexample.
    When $\texttt{Fixed = TRUE}$, TLC proves that \texttt{NoInformationLeak} holds across all possible interleavings.
\end{itemize}


\end{document}
